% sections/stress_test.tex
\section{Treadmill Stress Test (Bruce Protocol)}

The patient underwent a symptom-limited graded exercise stress test using the standard Bruce Protocol to evaluate for inducible ischemia and chronotropic response following clinical guidelines~\cite{gibbons2002acc}.

\subsection{Stress Test Performance Summary}

\begin{table}[H]
\centering
\caption{Exercise Stress Test Stage Parameters}
\label{tab:stress_stages}
\begin{tabularx}{\textwidth}{l c c c c c X}
\toprule
\textbf{Stage} & \textbf{Time (min)} & \textbf{Speed (mph)} & \textbf{Grade (\%)} & \textbf{METs} & \textbf{HR (bpm)} & \textbf{BP (mmHg)} \\
\midrule
Rest & 0:00 & -- & -- & 1.0 & 68 & 128/82 \\
Stage 1 & 3:00 & 1.7 & 10.0 & 4.6 & 95 & 138/84 \\
Stage 2 & 6:00 & 2.5 & 12.0 & 7.0 & 125 & 148/86 \\
Stage 3 & 9:00 & 3.4 & 14.0 & 10.1 & 152 & 158/88 \\
Stage 4 (Peak) & 10:30 & 4.2 & 16.0 & 12.9 & 165 & 166/90 \\
Recovery 1 & 11:30 & 1.7 (Walk) & 0.0 & 2.0 & 135 & 150/86 \\
Recovery 3 & 13:30 & -- & -- & 1.0 & 100 & 136/82 \\
Recovery 5 & 15:30 & -- & -- & 1.0 & 78 & 130/80 \\
\bottomrule
\end{tabularx}
\end{table}

\subsection{Heart Rate Response Curve}
The plot below illustrates the patient's heart rate trajectory during exercise and recovery compared to the age-predicted maximal heart rate (APMHR) and the target 85\% submaximal threshold.

\begin{figure}[H]
\centering
\begin{tikzpicture}
\begin{axis}[
    width=0.9\textwidth,
    height=7cm,
    xlabel={Elapsed Time (minutes)},
    ylabel={Heart Rate (beats per minute)},
    xmin=0, xmax=16,
    ymin=50, ymax=190,
    xtick={0, 3, 6, 9, 10.5, 11.5, 13.5, 15.5},
    xticklabels={Rest, Stg 1, Stg 2, Stg 3, Peak, Rec 1, Rec 3, Rec 5},
    ytick={60, 80, 100, 120, 140, 160, 180},
    grid=both,
    grid style={line width=.1pt, draw=gray!10},
    major grid style={line width=.2pt, draw=gray!25},
    legend pos=south west,
    legend cell align={left},
    axis line style={MedNavy, thick},
    tick style={thick, MedNavy}
]
    % Observed HR Curve
    \addplot[
        color=MedTeal,
        mark=*,
        thick,
        mark options={fill=white, draw=MedTeal, line width=1.5pt}
    ] coordinates {
        (0.0, 68)
        (3.0, 95)
        (6.0, 125)
        (9.0, 152)
        (10.5, 165)
        (11.5, 135)
        (13.5, 100)
        (15.5, 78)
    };
    \addlegendentry{Observed Heart Rate}

    % Target 85% HR Line (142 bpm)
    \addplot[
        color=MedNavy,
        dashed,
        thick
    ] coordinates {
        (0, 142) (16, 142)
    };
    \addlegendentry{85\% APMHR Target (142 bpm)}

    % 100% APMHR Line (167 bpm)
    \addplot[
        color=AlertRed,
        dotted,
        thick
    ] coordinates {
        (0, 167) (16, 167)
    };
    \addlegendentry{100\% APMHR Max (167 bpm)}
    
\end{axis}
\end{tikzpicture}
\caption{Heart Rate profile during exercise and recovery phases.}
\label{fig:hr_response}
\end{figure}

\subsection{Exercise Test Assessment}
\begin{itemize}
    \item \textbf{Exercise Capacity:} The patient completed 10 minutes and 30 seconds of exercise, terminating during Stage 4. Total workload achieved was 12.9 METs, indicating an \textbf{excellent} functional aerobic capacity for his age.
    \item \textbf{Hemodynamic Response:} Resting blood pressure was 128/82 mmHg, rising to a peak exercise blood pressure of 166/90 mmHg. Heart rate increased from a resting value of 68 bpm to a peak of 165 bpm, representing 99\% of his age-predicted maximal heart rate (APMHR = 167 bpm). Chronotropic index was normal. Blood pressure and heart rate recovery profiles were normal, with heart rate dropping by 30 bpm at 1 minute of recovery.
    \item \textbf{Symptomatology:} The test was terminated due to general fatigue and mild shortness of breath. No chest pain, pressure, or other anginal equivalents were reported during or after the test.
    \item \textbf{Electrocardiographic Response:} No significant baseline-corrected ST-segment deviations or T-wave changes were observed in any lead during exercise or recovery. No cardiac arrhythmias or ectopy were detected.
\end{itemize}
